
%\needspace{4\baselineskip}
\color{uniblue}{\subsection{\raggedright{Контроллер присоединения (КП)}}}
\color{black}

\begin{enumerate}[label=\arabic{section}.\arabic{subsection}.\arabic{enumi}, labelsep=4pt, leftmargin=0pt, itemindent=57pt, itemsep=0pt, parsep=5pt] % если что можно крутить itemsep

%%%%%%%%%%%%%%%%%%%%%%%%%%%%%%%%%%%%%%%%%%%%%%%%%%%%%%%%%%%%%%%%%%%%%%  ОБЩИЕ СВЕДЕНИЯ %%%%%%%%%%%%%%%%%%%%%%%%%%%%%%%%%%%%%%%%%%%%%%%%%%
\item Общие сведения

\begin{enumerate}[label=\arabic{section}.\arabic{subsection}.\arabic{enumi}.\arabic*, labelsep=4pt, leftmargin=0em, itemindent=65pt, parsep=0pt]

\item
Функциональный блок <<КП>> предназначен:
\begin{itemize}
\item для оперативного управления коммутационными аппаратами; 
\item для формирования сигнала положения коммутационного аппарата для целей АСУ ТП; 
\item для реализации оперативной блокировки управляемых коммутационных аппаратов.
\end{itemize}

\item
В функциональном блоке <<КП>> в составе устройства заложена возможность оперативного управления следующими коммутационными аппаратами:
\begin{itemize}
	\item силовой выключатель (В);
	\item заземляющий нож (ЗН);
	\item выкатной элемент 1 (ВЭ1);
	\item выкатной элемент 2 (ВЭ2).
\end{itemize}
\item
При выполнении коммутационной операции В, ЗН, или ВЭ информация об этом процессе будет отображаться общим выходным сигналом «КП: Идет переключение». Функциональная схема формирования общего сигнала представлена рисунке \ref{kp:img_perkl}.

\item
Ввод функционального блока <<КП>> в работу на этапе параметрирования устройства осуществляется программным переключателем SGF1 <<Ввод\_функции>>. Введенное состояние функционального блока соответствует активному состоянию сигнала <<КП: Ввод>>.
\item
Функциональный блок <<КП>> может быть выведен из работы оперативно путем активации сигнала <<ОВ КП>>, о чем сигнализирует активное состояние сигнала <<КП: Оперативный вывод>>.
\item
Логика управления режимами работы функционального блока <<КП>> выполнена в соответствии с рисунком \ref{kp:img_gen}. В таблице \ref{kp:tbl_gen} приведены параметры, необходимые для настройки функционального блока <<КП>>. 

\vspace{3mm}
\begin{figure}[H]
\centering

\includegraphics[width=1\textwidth,height=1\textheight,keepaspectratio]{\fbpath/312. КП ВВ/_latex/img/1.pdf}
\captionof{figure}{Функциональная схема управления режимами работы функционального блока <<КП>>}\label{kp:img_gen}
\end{figure}

\vspace{3mm}
\begin{figure}[H]
\centering
\includegraphics[width=0.6\textwidth,height=0.6\textheight,keepaspectratio]{\fbpath/312. КП ВВ/_latex/img/2.pdf}
\captionof{figure}{Функциональная схема формирования общего сигнала переключения коммутационного аппарата}\label{kp:img_perkl}
\end{figure}

\small
\begin{longtable}{|>{\centering\arraybackslash}m{5.3cm}|>{\centering\arraybackslash}m{3.3cm}|>{\centering\arraybackslash}m{4.2cm}|>{\centering\arraybackslash}m{1.8cm}|>{\centering\arraybackslash}m{1cm}|}
\caption{Параметры для настройки функционального блока <<КП>>\hfill\vspace{-0.5\baselineskip}}\label{kp:tbl_gen}\\ 
\hline
\rowcolor{gray!30}
Параметр (Параметр на ИЧМ) & Условное обозначение на схеме & Значение/ Диапазон & Единица измерения & Шаг \\ 
\hline
\endfirsthead
\caption*{\hspace{3pt}\emph{Продолжение таблицы \ref{kp:tbl_gen}\hfill\vspace{-0.5\baselineskip}}} \\ % сделано по ГОСТ 2.105 п.6.8.7
\hline
\rowcolor{gray!30}
Параметр (Параметр на ИЧМ) & Условное обозначение на схеме & Значение/ Диапазон & Единица измерения & Шаг \\ 
%\hline
\endhead
%\hline
\endfoot
%\hline
\endlastfoot
%==+t1*LLN0|VV_SWCTRL

%===t1
\end{longtable}
\normalsize

\end{enumerate}
%%%%%%%%%%%%%%%%%%%%%%%%%%%%%%%%%%%%%%%%%%%%%%%%%%%%%%%%%%%%%%%%%%%%%%  КОНЕЦ ОБЩИЕ СВЕДЕНИЯ %%%%%%%%%%%%%%%%%%%%%%%%%%%%%%%%%%%%%%%%%%%%%%%%%%

%%%%%%%%%%%%%%%%%%%%%%%%%%%%%%%%%%%%%%%%%%%%%%%%%%%%%%%%%%%%%%%%%%%%%%  УПРАВЛЕНИЕ ВЫКЛЮЧАТЕЛЕМ %%%%%%%%%%%%%%%%%%%%%%%%%%%%%%%%%%%%%%%%%%%%%%%%%%

\item Управление выключателем (УВ)

\begin{enumerate}[label=\arabic{section}.\arabic{subsection}.\arabic{enumi}.\arabic*, labelsep=4pt, leftmargin=0em, itemindent=65pt, parsep=0pt]

\item
Функция <<УВ>> формирует команды на отключение и включение выключателя при оперативном управлении.

\item
Местное управление коммутационным аппаратом осуществляется только посредством кнопок \tikz[baseline={(0,-0.00ex)}]{\draw[line width=1.6pt] (0,0) rectangle (1em,1em);\node[anchor=center] at (0.5em,0.5em) {\textbf{I}};}~–~<<Включить>> и \tikz[baseline={(0,-0.00ex)}]{\draw[line width=1.6pt] (0,0) rectangle (1em,1em);\node[anchor=center] at (0.5em,0.5em) {\textbf{O}};}~–~<<Отключить>>. При этом необходимо перевести устройство в режим <<Местное>> с помощью заданной функциональной кнопки. В данном случае управление по локальной сети и с помощью выносных кнопок или ключей будет невозможно.
\item
Дистанционное управление может реализоваться через дискретные входы, которые сконфигурированы как <<Включить В от ТУ>> и <<Отключить В от ТУ>>, а также по локальной сети. В дистанционном режиме управления команды по локальной сети предусмотрены постоянно.
\item
Управление выключателем также осуществляется с помощью пульта управления (ПУ), который предназначен для формирования команд управления выключателем путем автоматической подачи на катушки электромагнитов электрического напряжения. Управление с ПУ может реализоваться через дискретные входы, которые сконфигурированы как <<Включить В от ПУ>> и <<Отключить В от ПУ>>.

\item
Команда оперативного отключения <<УВ: Отключить>> формируется в следующих случаях:
\begin{itemize}
	\item при появлении команды на оперативное отключение выключателя (<<Отключить В~от~ПУ>> или <<Откл. В~ИЧМ>>) и выбранным управлением по месту (сигнал <<Местное>> находится в активном состоянии);
	\item при появлении команды на оперативное отключение выключателя (<<Отключить В~от~ТУ>> или <<Откл. В~АСУ>>) и выбранным дистанционным управлением (сигнал <<Дистанционное>> находится в активном состоянии).
\end{itemize}
\item
Прохождение команды на отключение блокируется при наличии следующих сигналов:
\begin{itemize}
	\item запрет отключения от функции контроля силового выключателя <<КСВ: Блокировка откл>>;
	\item отключенное положение силового выключателя <<В отключен (б/к)>> (с контролем сигнала <<В включен~(б/к)>>);
	\item превышение времени переключения силового выключателя <<УВ: Прев. времени~пер.>>;
	\item местное управление силовым выключателем (положение переключателя <<М/Д>> в приводе выключателя) <<Ключ М/Д привода В>>.
\end{itemize}
\item
В устройстве применяется режим включения выключателя с контролем синхронизма, который контролируется уставкой SGF2 <<Контроль\_синхр>>. Если уставка SGF2 введена в работу, то при наличии разрешающих условий для включения выключателя и активного сигнала разрешения оперативного включения из функции КС <<Разреш.опер.вкл.В>> формируется команда включения с контролем синхронизма.

\item
Команда оперативного включения <<УВ: Включить>> без контроля синхронизма формируется в следующих случаях:
\begin{itemize}
	\item при появлении команды на оперативное включение выключателя (<<Включить В~от~ПУ>> или <<Вкл. В~ИЧМ>>) и выбранным управлением по месту (сигнал <<Местное>> находится в активном состоянии);
	\item при появлении команды на оперативное включение выключателя (<<Включить В~от~ТУ>> или <<Вкл. В~АСУ>>) и выбранным дистанционным управлением (сигнал <<Дистанционное>> находится в активном состоянии).
\end{itemize}

\item
Прохождение команды на включение блокируется при наличии следующих сигналов:
\begin{itemize}
	\item запрет включения от функции контроля силового выключателя <<КСВ: Блокировка вкл>>;
	\item включенное положение силового выключателя <<В~включен (б/к)>> (с контролем сигнала <<В~отключен~(б/к)>>);
	\item превышение времени переключения силового выключателя <<УВ: Прев. времени~пер.>>;
	\item местное управление силовым выключателем (положение переключателя <<М/Д>> в приводе выключателя) <<Ключ М/Д привода В>>;
	\item появление сигнала аварийного отключения силового выключателя <<КСВ: В аварийно отключен>> (прохождение сигнала определяется состоянием программного переключателя SGF3 <<Бл\_вкл\_от\_авар\_откл>>).
\end{itemize}

\item  
При наличии введенной уставки включения выключателя с контролем синхронизма SGF2 и превышение времени регулируемой уставки <<T3>> формируется выходной сигнал задержки включения для сброса команды.

\item
При наличии сигналов <<УВ: Включить>> и <<УВ: Отключить>> длительностью больше, чем время уставки <<T1>> формируется сигнал превышения допустимого времени переключения -- <<УВ: Прев. времени пер.>>. Активное состояние операции включения или отключения сигнализируется выходным сигналом <<УВ: Идет переключение>>.

\item
Логика функции <<УВ>> формирует сигналы состояния выключателя <<УВ: Отключено>>, <<УВ: Включено>>, <<УВ: Не определено>> и <<УВ: Неисправность>>, при этом сигналы <<УВ: Неисправность>> и <<УВ: Не определено>> выдаются с выдержкой времени, задаваемой уставкой Т2 <<Тблк>>.

\item
Ввод функции <<УВ>> в работу на этапе параметрирования устройства осуществляется программным переключателем SGF1 <<Ввод\_функции>>. Информация о введенной функции отображается выходным сигналом <<УВ: Ввод>>.

\item
Алгоритм работы <<УВ>> выполнен в соответствии с рисунком \ref{kp:img_uv}. В таблице \ref{kp:tbl_uv} приведены параметры, необходимые для настройки функции <<УВ>>.

\vspace{3mm}
\begin{figure}[H]
\centering
\includegraphics[width=1\textwidth,height=1\textheight,keepaspectratio]{\fbpath/312. КП ВВ/_latex/img/uv.pdf}
\captionof{figure}{Функциональная схема <<УВ>>}
\label{kp:img_uv}
\end{figure}

\small
\begin{longtable}{|>{\centering\arraybackslash}m{5.3cm}|>{\centering\arraybackslash}m{3.3cm}|>{\centering\arraybackslash}m{4.2cm}|>{\centering\arraybackslash}m{1.8cm}|>{\centering\arraybackslash}m{1cm}|}
\caption{Параметры для настройки функции <<УВ>>\hfill\vspace{-0.5\baselineskip}}\label{kp:tbl_uv}\\ 
\hline
\rowcolor{gray!30}
Параметр (Параметр на ИЧМ) & Условное обозначение на схеме & Значение/ Диапазон & Единица измерения & Шаг \\ 
\hline
\endfirsthead
\caption*{\hspace{3pt}\emph{Продолжение таблицы \ref{kp:tbl_uv}\hfill\vspace{-0.5\baselineskip}}} \\ % сделано по ГОСТ 2.105 п.6.8.7
\hline
\rowcolor{gray!30}
Параметр (Параметр на ИЧМ) & Условное обозначение на схеме & Значение/ Диапазон & Единица измерения & Шаг \\ 
%\hline
\endhead
%\hline
\endfoot
%\hline
\endlastfoot
%==+t1*CBCSWI1|VV_SWCTRL
\centering Ввод функции в работу (Ввод\_функции) & \centering SGF1 & \centering 0 = Не предусмотрено\\1 = Предусмотрено & \centering -- & \centering \arraybackslash -- \\
\hline
\centering Включение с контролем синхронизма (Контроль\_синхр) & \centering SGF2 & \centering 0 = Не предусмотрено\\1 = Предусмотрено & \centering -- & \centering \arraybackslash -- \\
\hline
\centering Блокировка включения от аварийного отключения В (Бл\_вкл\_от\_авар\_откл) & \centering SGF3 & \centering 0 = Не предусмотрено\\1 = Предусмотрено & \centering -- & \centering \arraybackslash -- \\
\hline
\centering Допустимое время переключения выключателя (Tперекл) & \centering T1 & \centering 0,01 ... 10,00 & \centering с & \centering \arraybackslash 0,01 \\
\hline
\centering Выдержка времени подавления выдачи положения "Не определено" (Tблк) & \centering T2 & \centering 0,01 ... 10,00 & \centering с & \centering \arraybackslash 0,01 \\
\hline
\centering Длительность задержки включения с контролем синхронизма (Тзад\_вкл\_с\_КС) & \centering T3 & \centering 1,00 ... 60,00 & \centering с & \centering \arraybackslash 0,01 \\
\hline
%===t1
\end{longtable}
\normalsize

\end{enumerate}
%%%%%%%%%%%%%%%%%%%%%%%%%%%%%%%%%%%%%%%%%%%%%%%%%%%%%%%%%%%%%%%%%%%%%% КОНЕЦ УПРАВЛЕНИЕ ВЫКЛЮЧАТЕЛЕМ %%%%%%%%%%%%%%%%%%%%%%%%%%%%%%%%%%%%%%%%%%%%%%%%%%

%%%%%%%%%%%%%%%%%%%%%%%%%%%%%%%%%%%%%%%%%%%%%%%%%%%%%%%%%%%%%%%%%%%%%%  УПРАВЛЕНИЕ ЗАЗЕМЛЯЮЩИМ НОЖОМ %%%%%%%%%%%%%%%%%%%%%%%%%%%%%%%%%%%%%%%%%%%%%%%%%%
\item Управление заземляющим ножом (УЗН)

\begin{enumerate}[label=\arabic{section}.\arabic{subsection}.\arabic{enumi}.\arabic*, labelsep=4pt, leftmargin=0em, itemindent=65pt, parsep=0pt]

\item
Функция <<УЗН>> формирует команды на отключение и включение зазмляющего ножа при оперативном управлении.

\item
Команда оперативного отключения <<УЗН: Отключить>> формируется в следующих случаях:
\begin{itemize}
	\item при появлении команды на оперативное отключение заземляющего ножа (<<Откл. ЗН ИЧМ>>), выбранным управлением по месту (сигнал <<Местное>> находится в активном состоянии) и наличием разрешающего сигнала от схемы оперативной блокировки <<ОБР ЗН: Разрешить управ.>>;
	\item при появлении команды на оперативное отключение заземляющего ножа (<<Откл. ЗН АСУ>>), выбранным дистанционным управлением (сигнал <<Дистанционное>> находится в активном состоянии) и наличием разрешающего сигнала от схемы оперативной блокировки <<ОБР ЗН: Разрешить управ.>>.
\end{itemize}

\item
Прохождение команды на отключение блокируется при наличии следующих сигналов:
\begin{itemize}
	\item отключенное положение заземляющего ножа <<ЗН отключен (б/к)>> (с контролем сигнала <<ЗН включен~(б/к)>>);
	\item превышение времени переключения заземляющего ножа <<УЗН: Прев. времени~пер.>>.
\end{itemize}

\item
Команда оперативного включения <<УЗН: Включить>> формируется в следующих случаях:
\begin{itemize}
	\item при появлении команды на оперативное включение выключателя (<<Вкл. ЗН ИЧМ>>), выбранным управлением по месту (сигнал <<Местное>> находится в активном состоянии) и наличием разрешающего сигнала от схемы оперативной блокировки <<ОБР ЗН: Разрешить управ.>>;
	\item при появлении команды на оперативное включение выключателя (<<Вкл. ЗН АСУ>>), выбранным дистанционным управлением (сигнал <<Дистанционное>> находится в активном состоянии) и наличием разрешающего сигнала от схемы оперативной блокировки <<ОБР ЗН: Разрешить управ.>>.
\end{itemize}

\item
Прохождение команды на включение блокируется при наличии следующих сигналов:
\begin{itemize}
	\item включенное положение заземляющего ножа <<ЗН~включен (б/к)>> (с контролем сигнала <<ЗН~отключен~(б/к)>>);
	\item превышение времени переключения заземляющего ножа <<УЗН: Прев. времени~пер.>>.
\end{itemize}

\item
При наличии сигналов <<УЗН: Включить>> и <<УЗН: Отключить>> длительностью больше, чем время уставки <<T1>> формируется сигнал превышения допустимого времени переключения -- <<УЗН: Прев. времени пер.>>. Активное состояние операции включения или отключения сигнализируется выходным сигналом <<УЗН: Идет переключение>>.

\item
Логика функции <<УЗН>> формирует сигналы состояния выключателя <<УЗН: Отключено>>, <<УЗН: Включено>>, <<УЗН: Не определено>> и <<УЗН: Неисправность>>, при этом сигналы <<УЗН: Неисправность>> и <<УЗН: Не определено>> выдаются с выдержкой времени, задаваемой уставкой Т2 <<Тблк>>.

\item
Ввод функции <<УЗН>> в работу на этапе параметрирования устройства осуществляется программным переключателем SGF1 <<Ввод\_функции>>. Информация о введенной функции отображается выходным сигналом <<УЗН: Ввод>>.

\item
Алгоритм работы <<УЗН>> выполнен в соответствии с рисунком \ref{kp:img_uzn}. В таблице \ref{kp:tbl_uzn} приведены параметры, необходимые для настройки функции <<УЗН>>.

\vspace{3mm}
\begin{figure}[H]
\centering
\includegraphics[width=1\textwidth,height=1\textheight,keepaspectratio]{\fbpath/128. КП ОЛ/_latex/img/uzn.pdf}
\captionof{figure}{Функциональная схема <<УЗН>>}\label{kp:img_uzn}
\end{figure}

\small
\begin{longtable}{|>{\centering\arraybackslash}m{5.3cm}|>{\centering\arraybackslash}m{3.3cm}|>{\centering\arraybackslash}m{4.2cm}|>{\centering\arraybackslash}m{1.8cm}|>{\centering\arraybackslash}m{1cm}|}
\caption{Параметры для настройки функции <<УЗН>>\hfill\vspace{-0.5\baselineskip}}\label{kp:tbl_uzn}\\ 
\hline
\rowcolor{gray!30}
Параметр (Параметр на ИЧМ) & Условное обозначение на схеме & Значение/ Диапазон & Единица измерения & Шаг \\ 
\hline
\endfirsthead
\caption*{\hspace{3pt}\emph{Продолжение таблицы \ref{kp:tbl_uzn}\hfill\vspace{-0.5\baselineskip}}} \\ % сделано по ГОСТ 2.105 п.6.8.7
\hline
\rowcolor{gray!30}
Параметр (Параметр на ИЧМ) & Условное обозначение на схеме & Значение/ Диапазон & Единица измерения & Шаг \\ 
%\hline
\endhead
%\hline
\endfoot
%\hline
\endlastfoot
%==+t1*ESCSWI1|SWCTRL
\centering Ввод функции в работу (Ввод\_функции) & \centering SGF1 & \centering 0 = Не предусмотрено\\1 = Предусмотрено & \centering -- & \centering \arraybackslash -- \\
\hline
\centering Допустимое время переключения (Tперекл) & \centering T1 & \centering 0,05 ... 100,00 & \centering с & \centering \arraybackslash 0,01 \\
\hline
\centering Выдержка времени подавления выдачи положения "Не определено" (Тблк) & \centering T2 & \centering 0,05 ... 100,00 & \centering с & \centering \arraybackslash 0,01 \\
\hline
%===t1
\end{longtable}
\normalsize

\end{enumerate}

%%%%%%%%%%%%%%%%%%%%%%%%%%%%%%%%%%%%%%%%%%%%%%%%%%%%%%%%%%%% КОНЕЦ УПРАВЛЕНИЕ ЗАЗЕМЛЯЮЩИМ НОЖОМ %%%%%%%%%%%%%%%%%%%%%%%%%%%%%%%%%%%%%%%%%%%%%%%%%%

%%%%%%%%%%%%%%%%%%%%%%%%%%%%%%%%%%%%%%%%%%%%%%%%%%%%%%%%%%%%%%%%  УПРАВЛЕНИЕ ВЫКАТНЫМ ЭЛЕМЕНТОМ %%%%%%%%%%%%%%%%%%%%%%%%%%%%%%%%%%%%%%%%%%%%%%%%%%

\item Управление выкатным элементом (УВЭ)

\begin{enumerate}[label=\arabic{section}.\arabic{subsection}.\arabic{enumi}.\arabic*, labelsep=4pt, leftmargin=0em, itemindent=65pt, parsep=0pt]

\item
Функциональный блок <<КП>> позволяет управлять двумя выкатными элементами, для этого в нем предусмотрены две функции управления <<УВЭ1>> и <<УВЭ2>>. Функции <<УВЭ1>> и <<УВЭ2>> идентичны по реализации алгоритма управления и по задаваемым уставкам.

\item
Функция <<УВЭ>> формирует команды вкатить и выкатить выкатной элемент ячейки фидера при оперативном управлении.

\item
Команда операции выкатить ВЭ <<УВЭ: Выкатить>> формируется в следующих случаях:
\begin{itemize}
	\item при появлении команды выкатить от местного управления (<<Выкат. ВЭ ИЧМ>>), выбранным управлением по месту (сигнал <<Местное>> находится в активном состоянии) и наличием разрешающего сигнала от схемы оперативной блокировки <<ОБР ВЭ: Разрешить управ.>>;
	\item при появлении команды выкатить от дистанционного управления (<<Выкат. ВЭ АСУ>>), выбранным дистанционным управлением (сигнал <<Дистанционное>> находится в активном состоянии) и наличием разрешающего сигнала от схемы оперативной блокировки <<ОБР ВЭ: Разрешить управ.>>.
\end{itemize}

\item
Прохождение команды выкатить ВЭ блокируется при наличии следующих сигналов:
\begin{itemize}
	\item контрольное положение ВЭ;
	\item превышение времени операции вкатить/выкатить ВЭ <<УВЭ: Прев. времени~пер.>>.
\end{itemize}

\item
Команда операции вкатить ВЭ <<УВЭ: Вкатить>> формируется в следующих случаях:
\begin{itemize}
	\item при появлении команды вкатить от местного управления (<<Вкат. ВЭ ИЧМ>>), выбранным управлением по месту (сигнал <<Местное>> находится в активном состоянии) и наличием разрешающего сигнала от схемы оперативной блокировки <<ОБР ВЭ: Разрешить управ.>>;
	\item при появлении команды вкатить от дистанционного управления (<<Вкат. ВЭ АСУ>>), выбранным дистанционным управлением (сигнал <<Дистанционное>> находится в активном состоянии) и наличием разрешающего сигнала от схемы оперативной блокировки <<ОБР ВЭ: Разрешить управ.>>.
\end{itemize}

\item
Прохождение команды вкатить ВЭ блокируется при наличии следующих сигналов:
\begin{itemize}
	\item рабочее положение ВЭ;
	\item превышение времени операции вкатить/выкатить ВЭ <<УВЭ: Прев. времени~пер.>>.
\end{itemize}

\item
При наличии сигналов <<УВЭ: Вкатить>> и <<УВЭ: Выкатить>> длительностью больше, чем время уставки <<T1>> формируется сигнал превышения допустимого времени операции -- <<УВЭ: Прев. времени пер.>>. Активное состояние операции включения или отключения сигнализируется выходным сигналом <<УВЭ: Идет переключение>>.

\item
Логика функции <<УВЭ>> формирует сигналы состояния выключателя <<УВЭ: Выкачен>>, <<УВЭ: Вкачен>>, <<УВЭ: Не определено>> и <<УВЭ: Неисправность>>, при этом сигналы <<УВЭ: Неисправность>> и <<УВЭ: Не определено>> выдаются с выдержкой времени, задаваемой уставкой Т2 <<Тблк>>.

\item
Ввод функции <<УВЭ>> в работу на этапе параметрирования устройства осуществляется программным переключателем SGF1 <<Ввод\_функции>>. Информация о введенной функции отображается выходным сигналом <<УВЭ: Ввод>>.

\item
Алгоритм работы <<УВЭ1>> выполнен в соответствии с рисунком \ref{kp:img_uve}, алгоритм управления для функции <<УВЭ2>> идентичный. В таблице \ref{kp:tbl_uve} приведены параметры, необходимые для настройки функции <<УВЭ1>>, параметры настройки функции <<УВЭ2>> аналогичны.

\vspace{3mm}
\begin{figure}[H]
\centering
\includegraphics[width=1\textwidth,height=1\textheight,keepaspectratio]{\fbpath/312. КП ВВ/_latex/img/uve.pdf}
\captionof{figure}{Функциональная схема <<УВЭ1>>}\label{kp:img_uve}
\end{figure}

\small
\begin{longtable}{|>{\centering\arraybackslash}m{5.3cm}|>{\centering\arraybackslash}m{3.3cm}|>{\centering\arraybackslash}m{4.2cm}|>{\centering\arraybackslash}m{1.8cm}|>{\centering\arraybackslash}m{1cm}|}
\caption{Параметры для настройки функции <<УВЭ1>>\hfill\vspace{-0.5\baselineskip}}\label{kp:tbl_uve}\\ 
\hline
\rowcolor{gray!30}
Параметр (Параметр на ИЧМ) & Условное обозначение на схеме & Значение/ Диапазон & Единица измерения & Шаг \\ 
\hline
\endfirsthead
\caption*{\hspace{3pt}\emph{Продолжение таблицы \ref{kp:tbl_uve}\hfill\vspace{-0.5\baselineskip}}} \\ % сделано по ГОСТ 2.105 п.6.8.7
\hline
\rowcolor{gray!30}
Параметр (Параметр на ИЧМ) & Условное обозначение на схеме & Значение/ Диапазон & Единица измерения & Шаг \\ 
%\hline
\endhead
%\hline
\endfoot
%\hline
\endlastfoot
%==+t1*WECSWI1|SWCTRL
\centering Ввод функции в работу (Ввод\_функции) & \centering SGF1 & \centering 0 = Не предусмотрено\\1 = Предусмотрено & \centering -- & \centering \arraybackslash -- \\
\hline
\centering Допустимое время переключения (Tперекл) & \centering T1 & \centering 0,05 ... 100,00 & \centering с & \centering \arraybackslash 0,01 \\
\hline
\centering Выдержка времени подавления выдачи положения "Не определено" (Тблк) & \centering T2 & \centering 0,05 ... 100,00 & \centering с & \centering \arraybackslash 0,01 \\
\hline
%===t1
\end{longtable}
\normalsize

\end{enumerate}
%%%%%%%%%%%%%%%%%%%%%%%%%%%%%%%%%%%%%%%%%%%%%%%%%%%%%%%%%%%%%%% КОНЕЦ УПРАВЛЕНИЕ ВЫКАТНЫМ ЭЛЕМЕНТОМ %%%%%%%%%%%%%%%%%%%%%%%%%%%%%%%%%%%%%%%%%%%%%%%%%%


%%%%%%%%%%%%%%%%%%%%%%%%%%%%%%%%%%%%%%%%%%%%%%%%%%%%%%%%%%%%%%%%%%%%%%  ОПЕРАТИВНАЯ БЛОКИРОВКА ОБЩЕЕ %%%%%%%%%%%%%%%%%%%%%%%%%%%%%%%%%%%%%%%%%%%%%%%%%%
\item Оперативная блокировка КА (ОБР КА)

\begin{enumerate}[label=\arabic{section}.\arabic{subsection}.\arabic{enumi}.\arabic*, labelsep=4pt, leftmargin=0em, itemindent=65pt, parsep=0pt]

\item
В устройстве реализовал алгоритм, который позволяет контролировать положение КА и выдавать сигналы на разрешение управления выключателем (В), выкатным элементом выключателя (ВЭ1), заземляющим ножом (ЗН) и выкатным элементом ТН до ввода (ВЭ2).

\item
Поясняющая схема первичной сети для алгоритмов оперативной блокировки КА приведена на рисунке \ref{kp:img_obrka}.

\vspace{3mm}
\begin{figure}[H]
\centering
\includegraphics[width=0.5\textwidth,height=0.5\textheight,keepaspectratio]{\fbpath/312. КП ВВ/_latex/img/obrka.pdf}
\captionof{figure}{Поясняющая схема первичной сети для алгоритмов оперативной блокировки КА}\label{kp:img_obrka}
\end{figure}

\end{enumerate}
%%%%%%%%%%%%%%%%%%%%%%%%%%%%%%%%%%%%%%%%%%%%%%%%%%%%%%%%%%%% КОНЕЦ ОПЕРАТИВНАЯ БЛОКИРОВКА ОБЩЕЕ %%%%%%%%%%%%%%%%%%%%%%%%%%%%%%%%%%%%%%%%%%%%%%%%%%

%%%%%%%%%%%%%%%%%%%%%%%%%%%%%%%%%%%%%%%%%%%%%%%%%%%%%%  Оперативная блокировка выключателя (ОБР В) %%%%%%%%%%%%%%%%%%%%%%%%%%%%%%%%%%%%%%%%%%%%%%%%%%
\item Оперативная блокировка выключателя (ОБР В)

\begin{enumerate}[label=\arabic{section}.\arabic{subsection}.\arabic{enumi}.\arabic*, labelsep=4pt, leftmargin=0em, itemindent=65pt, parsep=0pt]

\item
Условием ввода в работу разрешения управления выключателем является наличие входного сигнала <<УВ: Ввод>>. Контроль условий блокировки управления выключателем осуществляется с помощью уставки SGF1. Если уставка SGF1 выставлена в положение <<Не предусмотрено>>, то управление выключателем осуществляется без учета всех блокировок.

\item
Если уставка SGF1 выставлена в положение <<Предусмотрено>>, то разрешение управления выключателем формируется при выполнении следующих условий:
\begin{itemize}
	\item отсутствие операции переключения коммутационных аппаратов (сигнал <<КП: Идет переключение>>);
	\item активное состояние сигнала <<ВЭ: Выкачен>> для управления выключателем в контрольном (ремонтном) положениях или одновременное наличие следующих сигналов для разрешения 	управления выключателем в рабочем положении;
	\item активное состояние сигналов: <<ВЭ: Вкачен>>, <<ЗН: Отключен>>, <<КА1 Отключен>>, <<КА2 Отключен>>, если введен контроль данных сигналов соответствующими программными переключателями.	
\end{itemize}
\item
Внешние дискретные входы назначаются в проекте непосредственно под схему объекта. Блокировка включения выключателя выполняется при несоблюдении какого-либо из условий, предусмотренных алгоритмом логики <<ОБР В>>.

\item
Алгоритм работы <<ОБР В>> выполнен в соответствии с рисунком \ref{kp:img_obrv}. В таблице \ref{kp:tbl_obrv} приведены параметры, необходимые для настройки функции <<ОБР В>>.

\vspace{3mm}
\begin{figure}[H]
\centering
\includegraphics[width=0.85\textwidth,height=0.85\textheight,keepaspectratio]{\fbpath/128. КП ОЛ/_latex/img/obrv.pdf}
\captionof{figure}{Функциональная схема <<ОБР В>>}\label{kp:img_obrv}
\end{figure}

\small
\begin{longtable}{|>{\centering\arraybackslash}m{5.3cm}|>{\centering\arraybackslash}m{3.3cm}|>{\centering\arraybackslash}m{4.2cm}|>{\centering\arraybackslash}m{1.8cm}|>{\centering\arraybackslash}m{1cm}|}
\caption{Параметры для настройки функции <<ОБР В>>\hfill\vspace{-0.5\baselineskip}}\label{kp:tbl_obrv}\\ 
\hline
\rowcolor{gray!30}
Параметр (Параметр на ИЧМ) & Условное обозначение на схеме & Значение/ Диапазон & Единица измерения & Шаг \\ 
\hline
\endfirsthead
\caption*{\hspace{3pt}\emph{Продолжение таблицы \ref{kp:tbl_obrv}\hfill\vspace{-0.5\baselineskip}}} \\ % сделано по ГОСТ 2.105 п.6.8.7
\hline
\rowcolor{gray!30}
Параметр (Параметр на ИЧМ) & Условное обозначение на схеме & Значение/ Диапазон & Единица измерения & Шаг \\ 
%\hline
\endhead
%\hline
\endfoot
%\hline
\endlastfoot
%==+t1*CBCILO1|SWCTRL
\centering Ввод ОБР В в работу (ОБР\_В) & \centering SGF1 & \centering 0 = Не предусмотрено\\1 = Предусмотрено & \centering -- & \centering \arraybackslash -- \\
\hline
\centering Контроль ЗН (Контр\_ЗН) & \centering SGF2 & \centering 0 = Не предусмотрено\\1 = Предусмотрено & \centering -- & \centering \arraybackslash -- \\
\hline
\centering Контроль КА1 (Контр\_KA1) & \centering SGF3 & \centering 0 = Не предусмотрено\\1 = Предусмотрено & \centering -- & \centering \arraybackslash -- \\
\hline
\centering Контроль КА2 (Контр\_KA2) & \centering SGF4 & \centering 0 = Не предусмотрено\\1 = Предусмотрено & \centering -- & \centering \arraybackslash -- \\
\hline
%===t1
\end{longtable}
\normalsize

\end{enumerate}
%%%%%%%%%%%%%%%%%%%%%%%%%%%%%%%%%%%%%%%%%%%%%%%%%%%%%%%%%%%%%%%%%%%%%% КОНЕЦ Оперативная блокировка выключателя (ОБР В) %%%%%%%%%%%%%%%%%%%%%%%%%%%%%%%%%%%%%%%%%%%%%%%%%%

%%%%%%%%%%%%%%%%%%%%%%%%%%%%%%%%%%%%%%%%%%%%%%%%%%%%%%%%%%%%%%%%%%%%%%  Оперативная блокировка выкатного элемента (ОБР ВЭ) %%%%%%%%%%%%%%%%%%%%%%%%%%%%%%%%%%%%%%%%%%%%%%%%%%
\item Оперативная блокировка выкатного элемента (ОБР ВЭ)

\begin{enumerate}[label=\arabic{section}.\arabic{subsection}.\arabic{enumi}.\arabic*, labelsep=4pt, leftmargin=0em, itemindent=65pt, parsep=0pt]

\item
Условие ввода в работу разрешения управления ВЭ является наличие входного сигнала <<УВЭ: Ввод>>. Контроль условий блокировки управления ВЭ осуществляется с помощью уставки SG1. Если уставка SG1 (см. таблицу \ref{kp:tbl_gen}) выставлена в положение <<Не предусмотрено>>, то управление ВЭ осуществляется без учета всех блокировок.
\item
Принудительное деблокирование управления выкатным элементом от оперативной блокировки осуществляется активацией входного сигнала от ключа <<Программ. деблок. ВЭ>>.
\item
Разрешение управления ВЭ формируется при выполнении необходимых условий:
\begin{itemize}
	\item отсутствие процесса переключения коммутационных аппаратов (сигнал <<КП: Идет переключение>>);
	\item активное состояние сигнала <<ВЭ: Вкачен>> для разрешения команды выкатить ВЭ или одновременное наличие следующих сигналов для разрешения команды вкатить ВЭ: <<ВЭ: Выкачен>>, <<ЗН: Отключен>>, <<КА1 Отключен>>, <<КА2 Отключен>>, если введен контроль данных сигналов соответствующими уставками.
\end{itemize}

\item
Внешние дискретные входы назначаются в проекте непосредственно под схему объекта. Блокировка управления ВЭ выполняется при несоблюдении какого-либо из условий, предусмотренных алгоритмом логики <<ОБР ВЭ>>.

\item
Алгоритм работы <<ОБР ВЭ>> выполнен в соответствии с рисунком \ref{kp:img_obrve}. В таблице \ref{kp:tbl_obrve} приведены параметры, необходимые для настройки функции <<ОБР ВЭ>>.

\vspace{3mm}
\begin{figure}[H]
\centering

\includegraphics[width=1\textwidth,height=1\textheight,keepaspectratio]{\fbpath/128. КП ОЛ/_latex/img/obrve.pdf}
\captionof{figure}{Функциональная схема <<ОБР ВЭ>>}\label{kp:img_obrve}
\end{figure}

\small
\begin{longtable}{|>{\centering\arraybackslash}m{5.3cm}|>{\centering\arraybackslash}m{3.3cm}|>{\centering\arraybackslash}m{4.2cm}|>{\centering\arraybackslash}m{1.8cm}|>{\centering\arraybackslash}m{1cm}|}
\caption{Параметры для настройки функции <<ОБР ВЭ>>\hfill\vspace{-0.5\baselineskip}}\label{kp:tbl_obrve}\\ 
\hline
\rowcolor{gray!30}
Параметр (Параметр на ИЧМ) & Условное обозначение на схеме & Значение/ Диапазон & Единица измерения & Шаг \\ 
\hline
\endfirsthead
\caption*{\hspace{3pt}\emph{Продолжение таблицы \ref{kp:tbl_obrve}\hfill\vspace{-0.5\baselineskip}}} \\ % сделано по ГОСТ 2.105 п.6.8.7
\hline
\rowcolor{gray!30}
Параметр (Параметр на ИЧМ) & Условное обозначение на схеме & Значение/ Диапазон & Единица измерения & Шаг \\ 
%\hline
\endhead
%\hline
\endfoot
%\hline
\endlastfoot
%==+t1*WECILO1|SWCTRL
\centering Контроль КА1 (Контр\_КА1) & \centering SGF1 & \centering 0 = Не предусмотрено\\1 = Предусмотрено & \centering -- & \centering \arraybackslash -- \\
\hline
\centering Контроль КА2 (Контр\_КА2) & \centering SGF2 & \centering 0 = Не предусмотрено\\1 = Предусмотрено & \centering -- & \centering \arraybackslash -- \\
\hline
%===t1
\end{longtable}
\normalsize

\end{enumerate}
%%%%%%%%%%%%%%%%%%%%%%%%%%%%%%%%%%%%%%%%%%%%%%%%%%%%%%%%%%%%%%%%%%%%%% КОНЕЦ Оперативная блокировка выкатного элемента (ОБР ВЭ) %%%%%%%%%%%%%%%%%%%%%%%%%%%%%%%%%%%%%%%%%%%%%%%%%%

%%%%%%%%%%%%%%%%%%%%%%%%%%%%%%%%%%%%%%%%%%%%%%%%%%%%%%%%%%%%%%%%%%%%%%  Оперативная блокировка заземляющего ножа (ОБР ЗН) %%%%%%%%%%%%%%%%%%%%%%%%%%%%%%%%%%%%%%%%%%%%%%%%%%
\item Оперативная блокировка заземляющего ножа (ОБР ЗН)

\begin{enumerate}[label=\arabic{section}.\arabic{subsection}.\arabic{enumi}.\arabic*, labelsep=4pt, leftmargin=0em, itemindent=65pt, parsep=0pt]

\item
Условие ввода в работу разрешения управления ЗН является наличие входного сигнала <<УЗН: Ввод>>. Контроль условий блокировки управления ВЭ осуществляется с помощью уставки SG1. Если уставка SG1 (см. таблицу \ref{kp:tbl_gen}) выставлена в положение <<Не предусмотрено>>, то управление ЗН осуществляется без учета всех блокировок.

\item
Принудительное деблокирование управления заземляющим ножом от оперативной блокировки осуществляется активацией входного сигнала от ключа <<Программ. деблок. ЗН>>.

\item
Разрешение управления ЗН формируется при выполнении необходимых условий:
\begin{itemize}
	\item отсутствие процесса переключения коммутационных аппаратов (сигнал <<КП: Идет переключение>>);
	\item активное состояние сигнала <<ЗН: Включен>> для разрешения команды отключения ЗН или одновременное наличие следующих сигналов для разрешения команды включения ЗН: <<ВЭ: Выкачен>>, <<ЗН: Отключен>>, <<КА3 Отключен>>, <<КА4 Отключен>>, если введен контроль данных сигналов соответствующими уставками.
\end{itemize}

\item
Внешние дискретные входы назначаются в проекте непосредственно под схему объекта. Блокировка управления ЗН выполняется при несоблюдении какого-либо из условий, предусмотренных алгоритмом логики <<ОБР ЗН>>.

\item
Алгоритм работы <<ОБР ЗН>> выполнен в соответствии с рисунком \ref{kp:img_obrzn}. В таблице \ref{kp:tbl_obrzn} приведены параметры, необходимые для настройки функции <<ОБР ЗН>>.

\vspace{3mm}
\begin{figure}[H]
\centering
\includegraphics[width=1\textwidth,height=1\textheight,keepaspectratio]{\fbpath/128. КП ОЛ/_latex/img/obrzn.pdf}
\captionof{figure}{Функциональная схема <<ОБР ЗН>>}\label{kp:img_obrzn}
\end{figure}

\small
\begin{longtable}{|>{\centering\arraybackslash}m{5.3cm}|>{\centering\arraybackslash}m{3.3cm}|>{\centering\arraybackslash}m{4.2cm}|>{\centering\arraybackslash}m{1.8cm}|>{\centering\arraybackslash}m{1cm}|}
\caption{Параметры для настройки функции <<ОБР ЗН>>\hfill\vspace{-0.5\baselineskip}}\label{kp:tbl_obrzn}\\ 
\hline
\rowcolor{gray!30}
Параметр (Параметр на ИЧМ) & Условное обозначение на схеме & Значение/ Диапазон & Единица измерения & Шаг \\ 
\hline
\endfirsthead
\caption*{\hspace{3pt}\emph{Продолжение таблицы \ref{kp:tbl_obrzn}\hfill\vspace{-0.5\baselineskip}}} \\ % сделано по ГОСТ 2.105 п.6.8.7
\hline
\rowcolor{gray!30}
Параметр (Параметр на ИЧМ) & Условное обозначение на схеме & Значение/ Диапазон & Единица измерения & Шаг \\ 
%\hline
\endhead
%\hline
\endfoot
%\hline
\endlastfoot
%==+t1*ESCILO1|SWCTRL
\centering Контроль KA3 (Контр\_KA3) & \centering SGF1 & \centering 0 = Не предусмотрено\\1 = Предусмотрено & \centering -- & \centering \arraybackslash -- \\
\hline
\centering Контроль KA4 (Контр\_KA4) & \centering SGF2 & \centering 0 = Не предусмотрено\\1 = Предусмотрено & \centering -- & \centering \arraybackslash -- \\
\hline
%===t1
\end{longtable}
\normalsize

\end{enumerate}
%%%%%%%%%%%%%%%%%%%%%%%%%%%%%%%%%%%%%%%%%%%%%%%%%%%%%%%%%%%%%%%%%%%%%% КОНЕЦ Оперативная блокировка заземляющего ножа (ОБР ЗН) %%%%%%%%%%%%%%%%%%%%%%%%%%%%%%%%%%%%%%%%%%%%%%%%%%

\end{enumerate}
